\subsection{Status of the SU(2) Thread as Structural Attractor}
  \label{subsec:SU2-thread-status}

  The preceding subsections establish that $SU(2)$ emerges as the unique continuous
  limit of the chain of admissible finite relational structures
  $Q_8 \hookrightarrow 2I \hookrightarrow \mathrm{LPS}_p \hookrightarrow SU(2)$.
  This subsection clarifies the precise sense in which $SU(2)$ is an attractor and not
  merely a convenient limit.

  \paragraph{Finite resolution versus continuum.}
    At any fixed relational resolution, the admissible structure is a finite group.
    The admissibility ratio
    \begin{equation}
      R_p \;=\; \frac{A_{\rho_2}^{\max}}{A_{\rho_1}^{\max}}
      \;=\; \sqrt{\frac{\lambda_{\rho_1}}{\lambda_{\rho_2}}} \;>\; 1
    \end{equation}
    quantifies how much more amplitude the spinorial sector can sustain relative to the
    abelian sectors.
    For $Q_8$ on its 6-regular Cayley graph, $R_p = \sqrt{4/3}\approx 1.155$.
    For $2I$, the hierarchy is more pronounced: $R_p = \sqrt{14/9}\approx 1.247$.
    For the Lubotzky-Phillips-Sarnak family $\mathrm{LPS}_p$, the ratio grows with $p$
    while the graph remains Ramanujan, so the spectral gap is maximally efficient.

    As the resolution scale increases and the Cayley graph approximates the Haar measure
    on $SU(2)\cong S^3$, the distinction between abelian and non-abelian sectors dissolves:
    the continuous group has a single connected component, and the admissibility ratio
    $R_p\to 1$.
    The filter disappears not because it is switched off, but because $S^3$ already
    saturates the bounded-flux condition uniformly across all modes.

  \paragraph{Attractor in the space of finite relational structures.}
    The penalized spectral capacity $C_\alpha(G,S)$, which discounts groups achieving
    large amplitudes through constructive alignment of sectors, systematically favors
    binary polyhedral groups over their $SO(3)$ quotients for all $\alpha>0$.
    Binary covers dominate precisely because the central $\mathbb{Z}_2$ they retain
    prevents constructive interference between spinorial and abelian sectors, realising a
    more uniform distribution of admissible amplitude over the eigenspace decomposition.

    $SU(2)$ is therefore an attractor in the following operational sense: among all finite
    relational structures compatible with admissibility and bounded flux, those closest to
    $SU(2)$ in spectral capacity achieve the highest $C_\alpha$ for $\alpha>0$, and the
    sequence converges to $SU(2)$ as the resolution is refined.
    This is a structural selection, not a symmetry postulate.
